\clearpage
\fancyhead[L]{\small JUNO.}
\hypertarget{native-263-lower}{}
\begin{center}
JUNO.

\rule{0.6\linewidth}{0.4pt}

Monatliche Correspondenz. Band X. S. 552–555. Dezember 1804.

\rule{0.6\linewidth}{0.4pt}
\end{center}

\noindent—\quad—\quad—\par

Leider war der Himmel hier seit dem 6. November ständig bedeckt. Deshalb konnten wir Juno nicht weiter beobachten und Dr. \textsc{Gauss} keine neuen Daten liefern, mit denen er seine bereits berechneten Elemente hätte verbessern können. Doch die beiden Beobachtungen vom 20. und 21. Okt., die wir ihm mitteilten und die um $1\frac{1}{2}$ Minuten von seinen II. Elementen abwichen, reichten aus, um diesen ebenso unermüdlichen wie scharfsinnigen Astronomen zu veranlassen, sofort neue III. Elemente für Juno zu berechnen. Er schickte sie uns mit folgenden Bemerkungen. „Seit meinem letzten Brief habe ich die folgenden neuen III. Elemente berechnet. Dabei halfen mir Ihre Beobachtungen der Juno vom 20. und 21. Okt., die Sie mir freundlicherweise mitgeteilt hatten und die bereits um $1\frac{1}{2}$ Minuten von den zweiten Elementen abwichen.

\begin{center}
\begingroup\fontsize{8}{10}\selectfont\setlength{\tabcolsep}{3pt}\renewcommand{\arraystretch}{1.1}
\begin{tabular}{lr}
Epoche 1804, 30. Sept., um $0^{\mathrm u}$ im Seeberger Meridian . . & $22^\circ\ 34'\ 48''$ \\
Tägliche mittlere Bewegung . . . . . . . . . . . & $842''75$ \\
Aphel . . . . . . . . . . . . . . . . & $233^\circ\ 56'\ 6''$ \\
Logarithmus der großen Halbachse . . . . . . . & $0.426699$ \\
Exzentrizität . . . . . . . . . . . . . . . . & $0.263182$ \\
Aufsteigender Knoten . . . . . . . . . . . . . & $171^\circ\ 0'\ 0''$ \\
Neigung der Bahn . . . . . . . . . . . . . . & $12\ 52\ 48$ \\
\end{tabular}\endgroup
\end{center}

„Ich habe diese Elemente mit sämtlichen Ihrer Beobachtungen verglichen und folgende Übereinstimmung gefunden.
\SourceLeaf{264}
\begin{center}
\begingroup\fontsize{8}{10}\selectfont\setlength{\tabcolsep}{3pt}\renewcommand{\arraystretch}{1.1}
\begin{tabular}{r|r|r|r|r|r}
1804 & \shortstack{Mittlere Zeit\\auf dem\\Seeberg} & \shortstack{Berechnete\\Rektaszension\\der $\JunoSymbol$} & \shortstack{Berechnete südliche\\Deklination\\der $\JunoSymbol$} & \multicolumn{2}{c}{\shortstack{Differenz\\R—B\\in Rekt.\qquad in Dekl.}}\\ \hline
13. Sept. & $12^{\mathrm u}\ 31^{\mathrm m}\ 59^{\mathrm s}$ & $0^\circ\ 44'\ 55''8$ & $1^\circ\ 52'\ 34''3$ & $-0''6$ & $-2''7$ \\
14 & $12\ 27\ 27$ & $\phantom{000}\ 35\ 43.2$ & $2\ 5\ 35.3$ & $+2.2$ & $-0.2$ \\
15 & $12\ 22\ 53$ & $\phantom{000}\ 26\ 20.8$ & $\phantom{0}\ 18\ 42.2$ & $+0.8$ & $+3.7$ \\
17 & $12\ 13\ 45$ & $\phantom{000}\ 7\ 10.6$ & $\phantom{0}\ 45\ 9.2$ & $+2.5$ & $-3.9$ \\
18 & $12\ 9\ 10$ & $359\ 57\ 24.5$ & $\phantom{0}\ 58\ 28.2$ & $+4.0$ & $-2.7$ \\
20 & $.\ .\ .\ .$ & $\phantom{000}\ 37\ 35.7$ & $3\ 25\ 32.7$ & $.\ .\ .\ .$ & $+9.2$ \\
23 & $11\ 46\ 11$ & $\phantom{000}\ 7\ 25.9$ & $4\ 5\ 22.9$ & $+11.0$ & $-11.6$ \\
27 & $11\ 27\ 46$ & $358\ 27\ 7.7$ & $\phantom{0}\ 58\ 29.8$ & $+7.0$ & $-27.4$ \\
28 & $11\ 23\ 11$ & $\phantom{000}\ 17\ 9.6$ & $5\ 11\ 35.1$ & $+3.2$ & $-8.8$ \\
30 & $11\ 14\ 1$ & $357\ 57\ 28.6$ & $\phantom{0}\ 37\ 28.2$ & $-1.9$ & $-15.5$ \\
2. Okt. & $11\ 4\ 52$ & $\phantom{000}\ 38\ 15.5$ & $6\ 2\ 50.8$ & $+0.5$ & $-8.7$ \\
4 & $10\ 55\ 46$ & $\phantom{000}\ 19\ 40.9$ & $\phantom{0}\ 27\ 36.7$ & $+1.4$ & $-2.0$ \\
5 & $10\ 51\ 15$ & $\phantom{000}\ 10\ 41.1$ & $\phantom{0}\ 39\ 43.8$ & $-0.4$ & $+3.9$ \\
6 & $10\ 46\ 44$ & $\phantom{000}\ 1\ 54.0$ & $\phantom{0}\ 51\ 39.0$ & $-2.3$ & $+6.8$ \\
10 & $10\ 28\ 50$ & $356\ 29\ 14.9$ & $7\ 37\ 10.5$ & $+0.8$ & $+19.4$ \\
12 & $10\ 20\ 0$ & $\phantom{000}\ 14\ 37.8$ & $7\ 58\ 30.8$ & $-5.7$ & $+4.3$ \\
20 & $9\ 45\ 34$ & $355\ 29\ 55.3$ & $9\ 12\ 58.8$ & $+3.4$ & $-4.6$ \\
21 & $9\ 41\ 23$ & $\phantom{000}\ 26\ 2.7$ & $\phantom{0}\ 20\ 45.8$ & $+2.6$ & $+2.6$ \\
\end{tabular}\endgroup
\end{center}

„Dr. \textsc{Maskelyne} war so freundlich, mir noch die folgenden drei Beobachtungen mitzuteilen.

\begin{center}
\begingroup\fontsize{9}{11}\selectfont\setlength{\tabcolsep}{3pt}\renewcommand{\arraystretch}{1.1}
\begin{tabular}{r|r|r|r}
1804 & \shortstack{Mittlere Zeit\\in\\Greenwich} & \shortstack{Scheinbare\\Rektaszension\\der $\JunoSymbol$} & \shortstack{Scheinbare südliche\\Deklination\\der $\JunoSymbol$}\\ \hline
5. Okt. & $10^{\mathrm u}\ 51^{\mathrm m}\ 6^{\mathrm s}$ & $357^\circ\ 10'\ 24''6$ & $6^\circ\ 40'\ 1''1$ \\
9 & $10\ 33\ 9$ & $356\ 36\ 46.9$ & $7\ 26\ 21.3$ \\
17 & $9\ 58\ 10$ & $355\ 43\ 45.6$ & $8\ 47\ 18.5$ \\
\end{tabular}\endgroup
\end{center}

„Die III. Elemente stimmen mit diesen und den früheren Beobachtungen wie folgt überein.

\begin{center}
\begingroup\fontsize{8}{10}\selectfont\setlength{\tabcolsep}{3pt}\renewcommand{\arraystretch}{1.1}
\begin{tabular}{r|r|r|r|r|r}
1804 & \shortstack{Mittlere Zeit\\in\\Greenwich} & \shortstack{Berechnete\\Rektaszension\\der $\JunoSymbol$} & \shortstack{Berechnete südliche\\Deklination\\der $\JunoSymbol$} & \multicolumn{2}{c}{\shortstack{Differenz\\R—B\\in Rekt.\qquad in Dekl.}}\\ \hline
25. Sept. & $11^{\mathrm u}\ 36^{\mathrm m}\ 50^{\mathrm s}$ & $358^\circ\ 46'\ 55''7$ & $4^\circ\ 32'\ 26''8$ & $+7''7$ & $-1''4$ \\
29 & $11\ 18\ 27$ & $358\ 6\ 58.7$ & $5\ 24\ 58.3$ & $+6.4$ & $+1.6$ \\
5. Okt. & $10\ 51\ 6$ & $357\ 10\ 24.7$ & $6\ 40\ 6.0$ & $+0.1$ & $+4.9$ \\
9 & $10\ 33\ 9$ & $356\ 36\ 46.4$ & $7\ 26\ 28.0$ & $-0.5$ & $+6.7$ \\
17 & $9\ 58\ 10$ & $355\ 43\ 46.0$ & $8\ 47\ 27.1$ & $+0.4$ & $+8.6$ \\
\end{tabular}\endgroup
\end{center}

„Nach den III. Elementen ergibt sich für den künftigen Lauf der \emph{Juno} Folgendes. [Vom 2. Nov. 1804 bis zum 4. Jan. 1805 in Abständen von drei Tagen, Rektaszension, Deklination und Helligkeit.]

„Als Einheit wurde die Helligkeit gewählt, die der Planet in der Entfernung 1 von der Sonne und von der Erde hätte. Auf derselben Skala betrug sie

\begin{center}
\begingroup\fontsize{9}{11}\selectfont\setlength{\tabcolsep}{3pt}\renewcommand{\arraystretch}{1.1}
\begin{tabular}{lrl}
am 5. Sept. & $0.1378$ &  \\
— 12. — & $0.1484$ &  \\
— 3. Okt. & $0.1640$ & am größten in diesem Jahr \\
— 21. — & $0.1543$ &  \\
\end{tabular}\endgroup
\end{center}

„Es wird interessant sein zu sehen, wie lange \emph{Juno} diesmal sichtbar bleibt. Bei der nächsten Opposition Anfang März 1806 im Sternbild Löwe erreicht sie nur ein Viertel ihrer größten Helligkeit in diesem Jahr. Für den 31. Dezember 1805 ergibt sich für mich folgende Position
\SourceLeaf{265}
\begin{center}
Rekt. $176^\circ\ 45'$\quad Südl. Dekl. $2^\circ\ 44'$\quad Helligkeit $0.0284$.
\end{center}

\Needspace{17\baselineskip}
\indent Hier sind sämtliche Beobachtungen, die wir erhalten konnten.

\begin{center}
\begingroup\fontsize{9}{11}\selectfont\setlength{\tabcolsep}{3pt}\renewcommand{\arraystretch}{1.1}
\begin{tabular}{r|r|r|r}
1804 & \shortstack{Mittlere Zeit\\auf dem\\Seeberg} & \shortstack{Scheinbare\\Rektaszension\\der $\JunoSymbol$} & \shortstack{Scheinbare südliche\\Deklination\\der $\JunoSymbol$}\\ \hline
23. Okt. & $9^{\mathrm u}\ 33^{\mathrm m}\ 4^{\mathrm s}338$ & $355^\circ\ 19'\ 23''69$ & $9^\circ\ 35'\ 53''6$ \\
24 & $9\ 28\ 57.709$ & $355\ 16\ 42.51$ & $.\ .\ .\ .$ \\
30 & $9\ 4\ 54.644$ & $355\ 9\ 47.25$ & $10\ 18\ 43.8$ \\
5. Nov. & $8\ 41\ 56.099$ & $355\ 19\ 2.57$ & $.\ .\ .\ .$ \\
18 & $7\ 55\ 51.375$ & $356\ 34\ 46.58$ & $10\ 58\ 16.2$ \\
\end{tabular}\endgroup
\end{center}

\hfill[\textsc{von Zach.}]
